% Part: incompleteness
% Chapter: representability-in-q
% Section: undecidability

\documentclass[../../../include/open-logic-section]{subfiles}

\begin{document}

\olfileid{inc}{req}{und}
\olsection{అనిర్ణయనీయత}

సిద్ధాంతం~$\Th{T}$ భాషలోని ఏ వాక్యం~$!A$ ఇచ్చినా, ``$\Th{T}$ అనేది~$!A$ను
నిరూపిస్తుందా?'' అనే ప్రశ్నకు పరిమిత మెట్ల తరువాత తప్పక సరైన జవాబు ఇచ్చే
గణనా ప్రక్రియ ఏదీ లేకపోతే $\Th{T}$ను \emph{అనిర్ణయనీయమైనది} అంటాం. కాబట్టి
అంకగణిత భాషలో ఒక వాక్యం~$!A$ ఇచ్చినప్పుడు $\Th{Q}\Proves !A$ అవుతుందో లేదో
నిర్ణయించే గణనా ప్రక్రియ ఉంటే, అప్పుడు మాత్రమే $\Th{Q}$ నిర్ణయించదగినది.
దీన్ని మరింత ఖచ్చితంగా ఇలా అడగవచ్చు: $y$ అనేది~$\Th{Q}$లో నిరూపించదగిన
వాక్యపు గ్యోడెల్ సంఖ్య~$y$ అయినప్పుడు, అప్పుడు మాత్రమే నెరవేరే
$\Prov[\Th{Q}](y)$ సంబంధం పునరావృత్తమా? జవాబు: కాదు.

\begin{thm}
$\Th{Q}$ అనిర్ణయనీయమైనది; అంటే
\[
\Prov[\Th{Q}](y) \defiff \fn{Sent}(y) \land
\lexists[x][\Prf[\Th{Q}](x, y)]
\]
సంబంధం పునరావృత్తం కాదు.
\end{thm}

\begin{proof}
అది పునరావృత్తమని అనుకుందాం. అప్పుడు ఆగే సమస్యను ఇలా పరిష్కరించవచ్చు.
$e$, $n$ ఇచ్చినప్పుడు, $T$ అనేది
\olref[cmp][rec][nft]{thm:kleene-nf}లోని క్లీనీ విధేయమైతే, ఏదో~$s$కు
$T(e,n,s)$ అయినప్పుడు, అప్పుడు మాత్రమే $\cfind{e}(n)\fdefined$ అని తెలుసు.
$T$ ఆదిమ పునరావృత్త సంబంధం కనుక, $\Th{Q}$లో దానికి ఒక సూత్రం~$!B_T$
ప్రాతినిధ్యం చేస్తుంది; అంటే $T(e,n,s)$ అయినప్పుడు, అప్పుడు మాత్రమే
$\Th{Q}\Proves !B_T(\num{e},\num{n},\num{s})$. $\Th{Q}\Proves
!B_T(\num{e},\num{n},\num{s})$ అయితే $\Th{Q}\Proves
\lexists[y][!B_T(\num{e},\num{n},y)]$ కూడా. అలాంటి~$s$ ఏదీ లేకపోతే ప్రతి~$s$కూ
$\Th{Q}\Proves\lnot !B_T(\num{e},\num{n},\num{s})$. కానీ $\Th{Q}$
$\omega$-అవైరుధ్యమైనది; అంటే ప్రతి $n\in\Nat$కూ
$\Th{Q}\Proves\lnot !A(\num{n})$ అయితే
$\Th{Q}\Proves/\lexists[y][!A(y)]$. $\Th{Q}$ స్వీకృతాలు ప్రామాణిక
నమూనా~$\Struct{N}$లో సత్యమైనవి కనుక ఇది మనకు తెలుసు. అందువల్ల
$\Th{Q}\Proves/\lexists[y][!B_T(\num{e},\num{n},y)]$. మరో మాటలో, ఏదో~$s$కు
$T(e,n,s)$ అయినప్పుడు, అంటే $\cfind{e}(n)\fdefined$ అయినప్పుడు, అప్పుడు మాత్రమే
$\Th{Q}\Proves\lexists[y][!B_T(\num{e},\num{n},y)]$.
$e$, $n$ల నుంచి $\Gn{\lexists[y][!B_T(\num{e},
    \num{n}, y)]}$ను గణించవచ్చు; అలా చేసే ఆదిమ పునరావృత్త ప్రమేయాన్ని
$g(e,n)$గా తీసుకుందాం. కాబట్టి
\[
h(e, n) =
\begin{cases}
1 & \text{$\Prov[\Th{Q}](g(e, n))$ అయితే}\\
0 & \text{లేకపోతే}.
\end{cases}
\]
$\Prov[\Th{Q}]$ పునరావృత్తమైతే $h$ కూడా పునరావృత్తమని ఇది చూపుతుంది.
కానీ \olref[cmp][rec][hlt]{thm:halting-problem} వల్ల $h$ పునరావృత్తం కాదు.
కాబట్టి $\Prov[\Th{Q}]$ కూడా పునరావృత్తం కాలేదు.
\end{proof}

\begin{cor}
మొదటిస్థాయి తర్కం అనిర్ణయనీయమైనది.
\end{cor}

\begin{proof}
మొదటిస్థాయి తర్కం నిర్ణయించదగినదైతే, $\Th{Q}$లో నిరూప్యత కూడా
నిర్ణయించదగినదయ్యేది. ఎందుకంటే $!T$ అనేది $\Th{Q}$ స్వీకృతాల సంధానం అయితే,
$\Th{Q}\Proves !A$ అయినప్పుడు, అప్పుడు మాత్రమే $\Proves !T\lif !A$.
\end{proof}

\end{document}
